% Clean CV/Resume Template, by Bennet B <dev@bennet.cc>
% CC0, Public-Domain
% 
% Permission is hereby granted, free of charge, to any person obtaining a copy
% of this template and associated files (the "Template"), to deal
% in the Template without restriction, including without limitation the rights
% to use, copy, modify, merge, publish, distribute, sublicense, and/or sell
% copies of the Template, and to permit persons to whom the Template is
% furnished to do so, subject to the following conditions:
%
% The above copyright notice and this permission notice shall be included in all
% copies or substantial portions of the Template.
%
% THE TEMPLATE IS PROVIDED "AS IS", WITHOUT WARRANTY OF ANY KIND, EXPRESS OR
% IMPLIED, INCLUDING BUT NOT LIMITED TO THE WARRANTIES OF MERCHANTABILITY,
% FITNESS FOR A PARTICULAR PURPOSE AND NONINFRINGEMENT. IN NO EVENT SHALL THE
% AUTHORS OR COPYRIGHT HOLDERS BE LIABLE FOR ANY CLAIM, DAMAGES OR OTHER
% LIABILITY, WHETHER IN AN ACTION OF CONTRACT, TORT OR OTHERWISE, ARISING FROM,
% OUT OF OR IN CONNECTION WITH THE TEMPLATE OR THE USE OR OTHER DEALINGS IN THE
% TEMPLATE.
%
% based on Modern CV by Habib Semouma
% https://www.overleaf.com/latex/templates/modern-cv-slash-resume-template/vjrqdkpjckwj
%
% 
\documentclass[oneside]{article}

\usepackage{wallpaper}
\usepackage{geometry}
\usepackage[
    unicode=true,
    bookmarks=true,
    bookmarksnumbered=false,
    bookmarksopen=true,
    bookmarksopenlevel=1,
    breaklinks=false,
    pdfborder={0 0 0},
    backref=false,
    colorlinks=false
    ]{hyperref}
\usepackage{lastpage}
\usepackage[none]{hyphenat}
\usepackage{hyphsubst}
\usepackage{tabularx}
\usepackage{moresize}
\usepackage[document]{ragged2e}
% \usepackage{parskip}

\usepackage[scaled]{helvet}
\usepackage{fontawesome5}
\usepackage[defaultfam,tabular,oldstyle]{montserrat}
\usepackage[utf8]{inputenc}
\usepackage[T1]{fontenc}
\renewcommand*\oldstylenums[1]{{\fontfamily{Montserrat-TOsF}\selectfont #1}}

\usepackage{titlesec}
\usepackage{xcolor}
\usepackage{tikz}

\setlength{\parindent}{0pt}
\titleformat{\section}{\normalfont}{}{0pt}{}

\renewcommand{\arraystretch}{1.4}

\setlength\fboxrule{0pt}
\setlength\fboxsep{5pt}
% \setlength{\parskip}{.5\baselineskip plus 2pt}
% \renewcommand{\baselinestretch}{1.1}
\renewcommand{\baselinestretch}{0.9}

\titlespacing{\section}{0pt}{0.5ex plus .1ex minus .2ex}{0.2pc}

\newcolumntype{Y}{>{\RaggedRight\arraybackslash}X}

% Change PDF Meta Info here
\hypersetup{
    pdftitle={Julien Pontoire - CV - Français},
    pdfauthor={Julien Pontoire},
    pdfsubject={CV}
}

% Paper size
\geometry{
    a4paper,
    left=0pt,
    right=0pt,
    top=0pt,
    bottom=0pt,
    nohead,
    % includefoot,
    nomarginpar
}

% Background Color of the Sidebar Column
\definecolor{sidebg}{cmyk}{1, 0.02, 0, 0.56}
% Background Color of the Main Column
\definecolor{mainbg}{cmyk}{0, 0, 0.07, 0.04}

% Text Color of the Main Column
\definecolor{maintext}{cmyk}{1, 0.02, 0, 0.8}
% Text Color of the Sidebar Column
\definecolor{sidetext}{cmyk}{0, 0, 0.07, 0.04}

\pagecolor{mainbg}

\begin{document}
\setlength{\topskip}{0pt}\setlength{\footskip}{0pt}%
\fcolorbox{red}{sidebg}{%
    \begin{minipage}[t][\textheight-2\fboxsep-2\fboxrule][t]{\dimexpr0.3\textwidth-2\fboxrule-2\fboxsep\relax}
        \color{sidetext}
        %%%%%%%%%%%%%%%%%%%%%%%%%%%%%%%%%%%%%%%%%%%%%%%%%%%%
        % DESCRIPTION
        {\bfseries\scshape\HUGE Julien} \\
        {\bfseries\scshape\Huge Pontoire}
        \vspace{.3cm} \\
        Étudiant Ingénieur Informatique - Intelligence Artificielle et Sciences des Données
        
        \vspace{8pt} \\
        \rule{\linewidth}{0.4pt}
        \vspace{8pt} \\        
        %%%%%%%%%%%%%%%%%%%%%%%%%%%%%%%%%%%%%%%%%%%%%%%%%%%%
        % PERSONAL INFROMATION
        \phantomsection{}
        \addcontentsline{toc}{section}{Informations personnelles}
        \section*{\large Informations personnelles}
        \vspace{.2cm}
        \begin{tabularx}{\textwidth}{cY}
            \faPhone{}      & +33 7 82 75 34 71 \\[0.3cm]
            \faEnvelope{}   & \href{mailto:julien.pontoire@etu.utc.fr}{julien.pontoire@etu.utc.fr} \\[0.3cm]
            \faMapMarker{}  & France, Canada
        \end{tabularx}
        
        \vspace{8pt}
        \rule{\linewidth}{0.4pt}
        \vspace{8pt}
        
        %%%%%%%%%%%%%%%%%%%%%%%%%%%%%%%%%%%%%%%%%%%%%%%%%%%%%%%%%
        % LINKS
        \phantomsection{}
        \addcontentsline{toc}{section}{Liens}
        \section*{\large Liens}
        \vspace{.2cm}
        \begin{tabular}{cl}
        %    \faCode{}     & \href{https://example.com}{example.com}
            \faGithub{}   & \href{https://github.com/jpontoire}{github.com/jpontoire} \\[0.3cm]
            \faLinkedin{} & \href{https://www.linkedin.com/in/julien-pontoire-0a50a42b1/}{Julien Pontoire}
        \end{tabular}
        
        \vspace{8pt}
        \rule{\linewidth}{0.4pt}
        \vspace{8pt}

        %%%%%%%%%%%%%%%%%%%%%%%%%%%%%%%%%%%%%%%%%%%%%%%%%%%%%%%%%%%
        \phantomsection{}
        \addcontentsline{toc}{section}{Engagements}
        \section*{\large Engagements}
        \vspace{.2cm}
        \begin{tabularx}{\textwidth}{cY}
            & Élu étudiant au Conseil des Études et de la Vie Universitaire (2023 - 2025)
        \end{tabularx}
        
        \vspace{8pt}
        \rule{\linewidth}{0.4pt}
        \vspace{8pt}
        
        %%%%%%%%%%%%%%%%%%%%%%%%%%%%%%%%%%%%%%%%%%%%%%%%%%%%%%%%%%%%
        % YOUR SKILLS
        \phantomsection{}
        \addcontentsline{toc}{section}{Compétences}
        \section*{\large Compétences}
        \vspace{.2cm}
        \begin{tabularx}{\textwidth}{cY}
            \faCode{}        & C, Python, Rust, Java, Javascript, R, LISP \\[0.5cm]
            \faCubes{}        & Numpy, Pandas, SkLearn, PySpark, Ollama, BeautifulSoup, HTML, CSS, PHP, VueJS, Express, SpringBoot, React, PSQL \\[0.5cm]
            \faCogs{}        & Linux, Windows \\[0.5cm]
            \faLaptopCode{}  & Visual Studio Code, JetBrains, RStudio \\[0.5cm]
            \faGamepad{}     & Unreal Engine 5, Unity \\[0.5cm]
            \faToolbox{}     & Git, Perforce, UML, Phpmyadmin
        \end{tabularx}
        
        \vspace{8pt}
        \rule{\linewidth}{0.4pt}
        \vspace{8pt}
        
        %%%%%%%%%%%%%%%%%%%%%%%%%%%%%%%%%%%%%%%%%%%%
        % LANGUAGES
        \phantomsection{}
        \addcontentsline{toc}{section}{Langues}
        \section*{\large Langues}
        \vspace{.2cm}
        \begin{tabularx}{\textwidth}{cY}
            \faLanguage{} & Français - Langue Maternelle \\[0.3cm]
            \faLanguage{} & Anglais - C1 \\[0.3cm]
            \faLanguage{} & Japonais - A1
        \end{tabularx}        
        
    \end{minipage}%
}%
\hfill
\fcolorbox{red}{mainbg}{%
    \begin{minipage}[t][\dimexpr\textheight-2\fboxrule-2\fboxsep\relax][t]{\dimexpr0.7\textwidth-2\fboxrule-2\fboxsep\relax}
        \sloppy
        \color{maintext}
        \phantomsection{}
        {\LARGE \textbf{Stage Ingénieur en Sciences des Données et Intelligence Artificielle}} \\
        \vspace{0.2cm}
        {\scshape\large De Février 2026 à Juillet 2026}
        \vspace{0.3cm}
        
        %%%%%%%%%%%%%%%%%%%%%%%%%%%%%%%%%%%%%%%%%%%%%%%%%%%%%%%%%%%
        % EDUCATION
        \phantomsection{}
        \addcontentsline{toc}{section}{Éducation}
        \section*{\makebox[\linewidth][l]{\scshape\Large Éducation} \rule{\linewidth}{0.2pt}}

        {\large \textbf{Semestre d'études à l'étranger}} \\
        {\large \textbf{Université du Québec à Chicoutimi (UQAC)}} \\
        {\scshape\fontseries{light}\selectfont\footnotesize Chicoutimi \qquad Aout 2025 \textendash{} Décembre 2025} \\
        {Intelligence d'affaire, Développement de jeux vidéo, Architecture parallèle, Métaheuristiques} \\[1.5ex]

        {\large \textbf{Études d'ingénieur en informatique}} \\
        {\large \textbf{Université de Technologie de Compiègne (UTC)}} \\
        {\scshape\fontseries{light}\selectfont\footnotesize Compiègne \qquad Septembre 2021 \textendash{} Attendu en 2026} \\
        {Diplôme d'Ingénieur \textendash{ GPA: 4.63/5}} \\[1.5ex]
        
        {\large \textbf{Baccalauréat Général}} \\
        {\scshape\fontseries{light}\selectfont\footnotesize Lycée Boucher de Perthes, Abbeville \qquad Juillet 2021} \\
        {Mention Très Bien \textendash{} Spécialités Mathématiques et Sciences de l'Ingénieur} \\[1.5ex]

        %%%%%%%%%%%%%%%%%%%%%%%%%%%%%%%%%%%%%%%%%%%%%%%%%%%%%%%%%%
        % EXPERIENCE
        \phantomsection{}
        \addcontentsline{toc}{section}{Expérience}
        \section*{\makebox[\linewidth][l]{\scshape\Large Expérience} \rule{\linewidth}{0.2pt}}

        {\large \textbf{Stage assistant ingénieur au sein du médialab SciencesPo}}\\
        {\scshape\fontseries{light}\selectfont\footnotesize Septembre 2024 \textendash{} Février 2025} 
        \begin{itemize}
            \setlength{\itemsep}{-2pt}
            \item Amélioration et développement d'outils de collecte et traitement de données développés par le laboratoire
            \item Implémentation d’algorithmes de stemming (Porter, Carry) dans une librairie Rust
            \item Création d'algorithmes de collecte de données sur le site Reddit
            \item Développement d'un bookmarklet pour récupérer les résultats d'une recherche sur différents moteurs de recherche
            \item Collecte de données sur des sites de presse pour des fins de recherche scientifique
        \end{itemize}

        \phantomsection
        \addcontentsline{toc}{section}{Projets}
        \section*{\makebox[\linewidth][l]{\scshape\Large Projets} \rule{\linewidth}{0.2pt}}

        {\large \textbf{Mise en place d'un narrateur en réalité virtuelle}} \\
        {\scshape\fontseries{light}\selectfont\footnotesize Février 2025 \textendash{} Juillet 2025} \\
        \begin{itemize}
            \setlength{\itemsep}{-2pt}
            \item Projet de recherche : développement d'un narrateur via IA générative dans une scène de réalité mixte dans le cadre de l'exposition du mémorial de l'internement et de la déportation à Compiègne.
        \end{itemize}

        {\large \textbf{Prédiction de parcours étudiants}} \\
        {\scshape\fontseries{light}\selectfont\footnotesize Février 2025 \textendash{} Juillet 2025} \\
        \begin{itemize}
            \setlength{\itemsep}{-2pt}
            \item {\looseness=-1 Analyse d’un jeu de données d’enseignement supérieur pour prédire le devenir académique des étudiants.}
            \item Application de méthodes supervisées (régression logistique, forêts aléatoires) après une ACP et exploration non supervisée.
        \end{itemize}


        {\large \textbf{Data Warehouse et Business Intelligence}} \\
        {\scshape\fontseries{light}\selectfont\footnotesize Février 2025 \textendash{} Juin 2025} \\
        \begin{itemize}
            \setlength{\itemsep}{-2pt}
            \item Conception d’un processus ETL (Python)
            \item Création de pipeline de données (Snowflake)
            \item Réalisation de tableaux de bord interactifs (Power BI)
        \end{itemize}
        
        {\large \textbf{Cartographie de Reddit avec Gephi}} \\
        {\scshape\fontseries{light}\selectfont\footnotesize Septembre 2023 \textendash{} Décembre 2023} \\
        \begin{itemize}
            \setlength{\itemsep}{-2pt}
            \item Utilisation de l'API de Reddit
            \item Filtrage des données avec Python
            \item Classification et analyse de clusters
            \item Mise en carte avec Gephi
        \end{itemize}
    \end{minipage}
}%
\end{document}